% =========================================================
\chapter{C03 --- Méthodes d'analyse : dipôles, Thévenin, Norton, superposition}
% =========================================================

\begin{rembox}[title=En bref]
\textbf{De quoi parle ce chapitre ?} Le but n'est pas de reproduire tout ce qui se passe
\emph{dans} un circuit, mais de voir son \textbf{influence sur son environnement}. Un circuit
vu de l'extérieur par deux connexions est un \textbf{dipôle} (une boîte noire). Les théorèmes
de \textbf{Thévenin} et de \textbf{Norton} le réduisent à une source et une résistance ; le
théorème de \textbf{superposition} analyse une source à la fois. Tout vient encore des lois
de Kirchhoff. \\
\textbf{Les trois objectifs} : le principe des modèles \og boîte noire\fg{} ; l'équivalence
Thévenin--Norton (\og le passage secret entre le monde des tensions et celui des
courants\fg{}) ; les théorèmes de simplification.
\end{rembox}

\section{Le dipôle}

\slide{c03-03}{Trois expériences pour savoir si une boîte noire est active ou passive.}

\begin{defnbox}[title={Dipôle, passif ou actif}]
Un \textbf{dipôle} est une \textbf{boîte noire} reliée à son environnement par \textbf{deux
connexions} (les \og broches\fg{}). Peu importe ce qu'elle contient : seul compte son
comportement vu de l'extérieur.
\begin{itemize}
  \item \textbf{Dipôle passif} : aucune source capable de produire de l'énergie et/ou
        d'amplifier un signal (par exemple une combinaison de résistances).
  \item \textbf{Dipôle actif} : peut produire de l'énergie et/ou amplifier un signal ; il
        peut aussi contenir des composants passifs.
\end{itemize}
\end{defnbox}

\begin{propbox}[title=Les trois mesures qui identifient la nature du dipôle]
\begin{enumerate}
  \item \textbf{Tension à vide} ($I = 0$) : on lit $U_1$ ;
  \item \textbf{courant de court-circuit} ($U = 0$) : on lit $I \neq 0$ ;
  \item \textbf{tension avec une résistance externe} : $I = f(R)$, on lit $U_2$.
\end{enumerate}
Pour un dipôle \textbf{passif}, la tension mesurée est \textbf{toujours nulle} (avec ou sans
résistance externe), et le courant aussi. Sur la diapositive, en bas : des dipôles actifs
(source de tension + résistance, source de courant + résistance) et des dipôles passifs
(résistances seules).
\end{propbox}

\section{Le tripôle et le quadripôle}

\slide[0.7]{c03-04}{Tripôle (peu utilisé) et quadripôle.}

\begin{defnbox}[title=Tripôle et quadripôle]
\begin{itemize}
  \item Le \textbf{transistor} (MOS, bipolaire) est un \textbf{tripôle} : pas très pratique
        à exploiter, peu utilisé comme modèle.
  \item Le \textbf{quadripôle} (deux entrées, deux sorties) sert à \textbf{transformer les
        grandeurs physiques $I$, $V$} : un signal d'entrée (courant ou tension) influence un
        signal de sortie (courant ou tension). Les transformations les plus simples sont des
        \textbf{atténuations (filtres)} ou des \textbf{amplifications}. On modélise facilement
        un tripôle sous forme de quadripôle.
\end{itemize}
\end{defnbox}

\section{Rappel des simplifications locales}

\slide{c03-05}{Permutations et fusions.}

\begin{propbox}[title=Ce qu'on sait faire localement]
\begin{itemize}
  \item \textbf{Permutation série} dans un dipôle : l'ordre des éléments en série ne change
        rien vu de l'extérieur (on peut déplacer $U$, $R_1$, $R_2$ le long de la branche).
  \item \textbf{Permutation parallèle} : idem pour des éléments en parallèle.
  \item \textbf{Fusion série} : sources de tension en série ($U_3 = U_1 + U_2$, en tenant
        compte des sens) ; résistances en série ($R_3 = R_1 + R_2$).
  \item \textbf{Fusion parallèle} : sources de courant en parallèle ($I_3 = I_1 + I_2$) ;
        résistances en parallèle ($R_3 = R_1/\!/R_2$).
\end{itemize}
Les permutations servent à \textbf{rapprocher} des éléments qu'on peut ensuite fusionner.
\end{propbox}

\section{Application : ce qu'on a le droit de fusionner}

\slide{c03-06}{Expérience 1 (bonne réduction) et expérience 2 (charge avalée : conceptuellement faux).}

\begin{exbox}[title={Expérience 1 : réduction correcte}]
Source $U$, trois résistances $R$ en série, puis la charge $R_L$. Les trois $R$ fusionnent en
$3R$ ; on perd les tensions intermédiaires, mais elles n'intéressent pas la charge. Diviseur
résistif :
\[ V_L = U\cdot\frac{R_L}{3R + R_L} \]
\end{exbox}

\begin{exbox}[title={Expérience 2 : \og pas la bonne méthode\fg{}}]
Quatre $R$ en série, la charge $R_L$ en parallèle sur la dernière. Si on fusionne la charge
avec la résistance d'en face ($R/\!/R_L$) :
\[ V_L = U\cdot\frac{R/\!/R_L}{3R + (R/\!/R_L)} \]
$V_L$ reste facile à trouver, mais la charge a \textbf{virtuellement disparu} : \textbf{plus
aucune indication sur $I$} qui la traverse. $R_L$ a été intégrée au montage qu'on veut
modéliser : \textbf{conceptuellement faux}.
\end{exbox}

\begin{warnbox}[title=Morale]
Une réduction entraîne forcément une \textbf{perte d'information}. Les simplifications ne
s'appliquent pas de manière \og anarchique\fg{} : \textbf{on ne fusionne jamais la charge
qu'on étudie}.
\end{warnbox}

\section{Théorèmes de simplification : l'idée}

\slide{c03-07}{L'expérience 2 réduite en modèle de Thévenin ou de Norton.}

\begin{exbox}[title={Expérience 2, A.N. : $R = R_L = 1$\,k$\Omega$, $U = 14$\,V}]
\[ R_{EQ} = \frac{R\,R_L}{R + R_L} = \frac R2 \qquad
   V_L = \frac{R_{EQ}}{R_{EQ} + 3R}\,U = \frac{R/2}{R/2 + 3R}\,U = \frac17\cdot 14 = 2\text{ V} \]
(résultat à comparer avec la suite.)
\end{exbox}

Le gros montage (4 résistances + 1 source) peut se réduire à \textbf{une source et une
résistance} : soit une source de tension $V_{TH}$ avec $R_{IN}$ en série (\textbf{modèle
Thévenin}), soit une source de courant $I_{NO}$ avec $R_{IN}$ en parallèle (\textbf{modèle
Norton}). Ce n'est plus une simple fusion de composants, mais une approche globale incluant
\textbf{toutes} les sources et composants. Comme on sait fusionner tensions, courants et
résistances, il est légitime de trouver à la fin une source et une résistance. Question :
quelles expériences (avec une résistance, un ampèremètre A, un voltmètre V) mettent en
évidence la résistance et la source de la boîte noire ?

\begin{goatbox}
Peu importe la machinerie à l'intérieur de la grange. Vu de la porte, tout ce qui compte
c'est \og combien ça pousse\fg{} (la source) et \og combien ça résiste\fg{} (la résistance).
Thévenin et Norton disent que n'importe quelle grange linéaire, vue depuis deux bornes, se
comporte comme une seule pile avec une seule résistance.
\end{goatbox}

